\begin{thebibliography}{99}
\bibitem{BPR} S. Basu, R. Pollack and M.-F. Roy,
\emph{Algorithms in Real Algebraic Geometry}, 2nd ed., Springer, 2006.
\bibitem{BB} I. Berkes and B. Borda,
Random walks on the circle and Diophantine approximation,
\emph{J. London Math. Soc.} \textbf{108} (2023), 409--440.
\bibitem{Cassels} J. W. S. Cassels,
\emph{An Introduction to Diophantine Approximation}, Cambridge Tracts in Mathematics and Mathematical Physics, vol. 45, Cambridge University Press, 1957.
\bibitem{DGLPS} J. Dick, T. Goda, G. Larcher, F. Pillichshammer and K. Suzuki,
On the quasi-uniformity properties of quasi-Monte Carlo point sets and sequences---Part I: Lattices and Kronecker sequences,
arXiv:2502.06202v2, version of 13 February 2026, Section 3.3.
\bibitem{GL} S. Graf and H. Luschgy,
\emph{Foundations of Quantization for Probability Distributions}, Lecture Notes in Mathematics, vol. 1730, Springer, 2000.
\bibitem{LPV} C. H. Lee, E. Pyne and S. Vadhan,
Fourier growth of regular branching programs,
\emph{APPROX/RANDOM 2022}, LIPIcs \textbf{245}, Art. 2, 2:1--2:21.
\bibitem{RSV} O. Reingold, T. Steinke and S. Vadhan,
Pseudorandomness for regular branching programs via Fourier analysis,
\emph{RANDOM 2013}, LNCS \textbf{8096}, 655--670.
\bibitem{SVW} T. Steinke, S. Vadhan and A. Wan,
Pseudorandomness and Fourier growth bounds for width 3 branching programs,
\emph{Theory Comput.} \textbf{13} (2017), Art. 12, 1--50.
\bibitem{Vidyasagar} M. Vidyasagar,
\emph{Hidden Markov Processes: Theory and Applications to Biology}, Princeton University Press, 2014.

\bibitem{BF04} L. Benvenuti and L. Farina,
A tutorial on the positive realization problem,
\emph{IEEE Trans. Automat. Control} \textbf{49} (2004), 651--664,
Theorems 2--3 and Example 4.
\bibitem{Ben22} L. Benvenuti,
Minimal positive realizations: A survey,
\emph{Automatica} \textbf{143} (2022), Art. 110422.
\bibitem{BPP15} B. Balle, P. Panangaden and D. Precup,
A canonical form for weighted automata and applications to approximate minimization,
arXiv:1501.06841v3, 2015, Theorems 2 and 11.
\bibitem{FGS10} L. Finesso, A. Grassi and P. Spreij,
Two-step nonnegative matrix factorization algorithm for the approximate realization of hidden Markov models,
arXiv:1007.3435v1, 2010, Sections III--IV.
\bibitem{Ohta20} Y. Ohta,
On the realization of hidden Markov models and tensor decomposition,
arXiv:2008.11487v1, 2020, Theorems 6--9 and Remark 10.
\bibitem{BL06} Y. Bugeaud and M. Laurent,
Exponents of Diophantine approximation,
arXiv:math/0611354v1, 2006, Definition 1 and Section 2.
\bibitem{Tzeng92} W.-G. Tzeng,
A polynomial-time algorithm for the equivalence of probabilistic automata,
\emph{SIAM J. Comput.} \textbf{21} (1992), 216--227,
doi:10.1137/0221017.
\bibitem{Shitov18} Y. Shitov,
A universality theorem for nonnegative matrix factorizations,
arXiv:1606.09068v2, version of 5 April 2018, Theorem 2 and Corollary 3.
\bibitem{GS19} N. Gillis and Y. Shitov,
Low-rank matrix approximation in the infinity norm,
\emph{Linear Algebra Appl.} \textbf{581} (2019), 367--382;
arXiv:1706.00078v1, Lemmas 1--2 and Theorem 3.
\bibitem{MSZ23} S. Morozov, M. Smirnov and N. Zamarashkin,
On the optimal rank-1 approximation of matrices in the Chebyshev norm,
\emph{Linear Algebra Appl.} \textbf{679} (2023), 4--29;
arXiv:2212.01438, Theorem 5.3.
\bibitem{BKVH07} S. Boyd, S.-J. Kim, L. Vandenberghe and A. Hassibi,
A tutorial on geometric programming,
\emph{Optim. Eng.} \textbf{8} (2007), 67--127, Sections 2 and 8.
\bibitem{BV04} S. Boyd and L. Vandenberghe,
\emph{Convex Optimization}, Cambridge University Press, 2004, Section 5.8.
\bibitem{KKKR17} T. Kahle, K. Kubjas, M. Kummer and Z. Rosen,
The geometry of rank-one tensor completion,
\emph{SIAM J. Appl. Algebra Geom.} \textbf{1} (2017), 200--221;
arXiv:1605.01678v2, Section 2.
\bibitem{Sch86} A. Schrijver,
\emph{Theory of Linear and Integer Programming}, Wiley, 1986, Chapter 10.

\bibitem{BF03} L. Benvenuti and L. Farina,
Minimal positive realizations: A survey of recent results and open problems,
\emph{Kybernetika} \textbf{39} (2003), 217--228.
\bibitem{BF99} L. Benvenuti and L. Farina,
An example of how positivity may force realizations of ``large'' dimension,
\emph{Systems Control Lett.} \textbf{36} (1999), 261--266.
\bibitem{Yannakakis} M. Yannakakis,
Expressing combinatorial optimization problems by linear programs,
\emph{J. Comput. System Sci.} \textbf{43} (1991), 441--466.
\bibitem{RogersShephard} C. A. Rogers and G. C. Shephard,
The difference body of a convex body,
\emph{Arch. Math.} \textbf{8} (1957), 220--233.
\bibitem{NTZ26} J. Nie, X. Tang and J. Zhou,
Robust completion for rank-1 tensors with noises,
\emph{J. Sci. Comput.} \textbf{107} (2026), Art. 100,
doi:10.1007/s10915-026-03306-8.


\bibitem{Folland16} G. B. Folland,
\emph{A Course in Abstract Harmonic Analysis}, second edition,
CRC Press, 2016, Chapter 2.
\bibitem{Putinar93} M. Putinar,
Positive polynomials on compact semi-algebraic sets,
\emph{Indiana Univ. Math. J.} \textbf{42} (1993), 969--984.
\bibitem{Lasserre01} J.-B. Lasserre,
Global optimization with polynomials and the problem of moments,
\emph{SIAM J. Optim.} \textbf{11} (2001), 796--817,
doi:10.1137/S1052623400366802.
\bibitem{DoreaHennet99} C. E. T. D\'orea and J. C. Hennet,
$(A,B)$-invariant polyhedral sets of linear discrete-time systems,
\emph{J. Optim. Theory Appl.} \textbf{103} (1999), 521--542,
doi:10.1023/A:1021727806358.
\bibitem{Mejstrik20} T. Mejstrik,
Improved invariant polytope algorithm and applications,
arXiv:1812.03080v3, version of 21 June 2020.
\bibitem{Rabin63} M. O. Rabin,
Probabilistic automata,
\emph{Information and Control} \textbf{6} (1963), 230--245,
doi:10.1016/S0019-9958(63)90290-0.
\end{thebibliography}
